\chapter{Theoretical Backround and Related Work}

\section{Collision Detection and Bounding Volumes}

For a scene containing $n$ objects, directly testing every possible pair requires
$n(n-1)/2$ ($\mathcal{O}(n^2)$) pairwise collision tests in the worst case.
Collision detection is commonly divided into two stages, broad phase and narrow phase
The broad phase aims to effeciently eliminate pairs of objects that cannot collide and produce a set of potential collision candidate pairs.
The remaining candidates are then processed in the narrow phase, which performs collision tests on the actual geometry to determine wheter a collision occurs\cite[p.14]{ericson2004real}.

One approach for the acceleration of the broad phase is the use of Bounding Volumes.
A bounding volume is a simple geometric shape completely inclosing an object.
If the bounding volumes of two objects do not interesect, neither will the objects inside of the bounding volume.
In this work axis-aligned bounding-boxes are used as bounding volumes.
Axis-asligned bounding boxes are rectangular bounding volumes, which edges are aligned with the coordinate axes, making overlap tests computational inexpensive. \cite[pp.~75--81]{ericson2004real} \cite[p.10]{vandenbergen1997aabb}.

\section{Bounding Volume Hierarchies}

A bounding Volume Hierarchy (BVH) organizes multiple bounding volumes in a hierarchial tree-structure. 
In a binary BVH each leaf node of the tree encloses an primitive object, while internal nodes include all the primitives in its child nodes. \cite[pp.~235--238]{ericson2004real}

During collision detecition, the hierarchy can be traversed by comparing the bounding volumes of two nodes. 
If two bounding volumes do not overlap, all primitives represented by the subtree corresponding two the bounding volume can be discrded.
In case they do overlap, the subtrees will recursively be traversed until leaf nodes are reached, where primitive tests are conducted \cite[pp.~253--260]{ericson2004real}. From this we can conclude that the structure of the tree influences how much of the subtrees can actually be discarded. 

An effectiuve hierarchy should allow large parts of the set to be sorted out early.
Properties such as tightness and simplicity of the bounding volumes, 

\section{BVH Construction Strategies}

Ericson distringuishes between three different approaches to BVH construction, these being top-down, bottom-up and incremental construction \cite[pp.~239--253]{ericson2004real}.
Although algorithms within these categories differ in the heuristics used, those categories can be used to describe the general way the hierarchy is generated.

\subsection{Top-Down Construction}

Top-down construction begins with aöö primitives assigned to the root node.
The primitive set is then divided into subsets, which are assigned as childeren of the root node. This is then done recursively until a subset contains only one primitive \cite[pp.~240--244]{ericson2004real}.

\subsection{Bottom-Up Construction}

Bottom-up construction proceeds in the opposite direction. Instead of splitting up subsets, in bottom-up subsets are created from merging leaf nodes ccontaining individual primitives. 
Nodes are repeatedly selected according to a clustering or merging criteria and vombined under a new parent node, until only the root node of the tree is left. \cite[pp.~245--250]{ericson2004real}

\subsection{Incremental Construction}

In incremental tree construction, 

\newpage